\documentclass[a4paper,11pt]{article}

\usepackage{revy}
\usepackage[utf8]{inputenc}
\usepackage[T1]{fontenc}
\usepackage[danish]{babel}
\usepackage{amsmath}
\usepackage{amssymb}   % Required for blackboard bold (\mathbb) and fraktur (\mathfrak)
\usepackage{mathrsfs}


\revyname{MatematikRevyen}
\revyyear{2026}
\version{1}
\eta{$3$ minutter}
\status{Faerdig}

\title{Fuldstændiggøre $\mathbb{Q}$
}
\author{Adam Sonne Hansen '23}
\melody{Saveus og Guldimund "Vil du noget"}

\begin{document}
\maketitle

\begin{roles}
\role{S}[Cauchyfølge 1]
\role{S}[Cauchyfølge 2]
\end{roles}

To Cauchyfølger vil gerne konvergere mod kvadratrod 2. Rækkefølgen på hvem der synger og sådan er lidt ligemeget.

\begin{song}
  \sings{}[Omkvæd]
  Åh, jeg ku'
  Konvergere, mod kvadratrod 2
  Åh, men det' ik' en del
  Af legemet, du og jeg vi lever i
  Hey, hey vil du
  Fuldstændiggøre, fuldstændiggøre $\mathbb{Q}$
  Hey, hey vil du
  Fuldstændiggøre, fuldstændiggøre $\mathbb{Q}$

Fuldstændiggøre $\mathbb{Q}$

  \sings{}[Vers 1]
Der er huller i linjen, mellem tallene
Tager du en kugle, og en mængde ser du de snitter tomt
Ka' ik' konvergere, drømmer om mere, savner en grænseværdi
Savner at være tæt, kom, lad os skabe noget
Irrationalitet
Hvis vi ku' definer'
Et legeme med lidt mer'
Et legeme helt komplet
Hvor rationale tal er tæt'

\sings{}[Omkvæd]
  Åh, jeg ku'
  Konvergere, mod kvadratrod 2
  Åh, men det' ik' en del
  Af legemet, du og jeg vi lever i
  Hey, hey vil du
  Fuldstændiggøre, fuldstændiggøre $\mathbb{Q}$
  Hey, hey vil du
  Fuldstændiggøre, fuldstændiggøre $\mathbb{Q}$

  Fuldstændiggøre $\mathbb{Q}$
  
\sings{}[Vers 2]
(Så) Vi konstruerer, et legeme uden et hul i
Starter med definitionen af en ny  absolutværdi
Én med en stærkere version af en trekantsulighed
Én som der relaterer, til primtallenes delelighed
Ta'r man en følge her, og la'r den konverger'
Ser vi en følge, af Cauchyfølger, har (en) grænse
Legemet af $p$-adiske tal er så pænt 
Forskel af termer, de bli'r ingenting 
Hvis og kun hvis, $(x_n)$ er Cauchy.
Gid det bare var såd'n for enhver and'n absolutværdi
Dens egenskaber de føder en mærklig Topologi
Cirklernes indre punkter de er lig med centrum indeni
Her er vi alle tætte, ogs' de irrationelle 
Mængder er komplette, på samme måde som reelle
Rod to er i $\mathbb{Q}_p$, lad os konvergere
For Hensels Lemma si'r
At det eksisterer

\sings{}[Omkvæd]
  Åh, jeg ku'
  Konvergere, mod kvadratrod 2
  Åh, det er nu en del
  Af legemet, du og jeg vi lever i
  Hey, vi har nu
  Fuldstændiggjort, fuldstændiggjort $\mathbb{Q}$
  Hey, vi har nu
  Fuldstændiggjort,, fuldstændiggjort, fuldstændiggjort $\mathbb{Q}$

\end{song}

\end{document}